% ============================
% MATLAB language + style inspirado en el editor de MATLAB.
% ============================

\RequirePackage{listings}
\RequirePackage{xcolor}
\RequirePackage{courier}

% Colores del editor de MATLAB.
\definecolor{mygreen}{RGB}{0,128,0}
\definecolor{mylilas}{RGB}{170,55,241}
\definecolor{matlabblue}{RGB}{0,0,255}

% Definición del lenguaje MATLAB (ajustado al esquema del snippet)
\lstdefinelanguage{Matlab}{
  sensitive=true,
  % Palabras clave básicas de MATLAB
  keywords={
    break,case,catch,classdef,continue,else,elseif,end,for,function,global,if,
    otherwise,parfor,persistent,return,spmd,switch,try,while,properties,methods
  },
  % Extra del snippet
  morekeywords={matlab2tikz},
  morekeywords=[2]{1},
  % Palabras reservadas azules e identificadores negros.
  keywordstyle=\color{matlabblue},
  keywordstyle=[2]\color{black},
  identifierstyle=\color{black},
  % Comentarios y cadenas
  comment=[l]\%,                      % comentario con %
  morecomment=[s]{\%\{}{\%\}},        % bloque %{ ... %}
  commentstyle=\color{mygreen},
  morestring=[b]',                    % strings con comilla simple
  stringstyle=\color{black},
  % Varios
  alsoletter={...}                    % para que '...' no rompa tokens
}

% Estilo listo para usar (equivale al \lstset{...} del snippet)
\lstdefinestyle{myMatlab}{
  language=Matlab,
  breaklines=true,
  showstringspaces=false,
  numbers=left,
  numberstyle={\tiny\color{black}},
  numbersep=9pt,
  basicstyle=\color{black}\ttfamily\small,
  frame=single,
  columns=fullflexible,
  keepspaces=true,
  tabsize=2,
  captionpos=b,
  % Las palabras enfatizadas también son palabras reservadas.
  emph={for,end,break},
  emphstyle=\color{matlabblue}
}

% --- Entorno inline: \begin{MatlabCode}[opciones] ... \end{MatlabCode}
\lstnewenvironment{MatlabCode}[1][]{%
  \lstset{style=myMatlab,#1}%
}{}

% --- Atajo para incluir archivos .m: \MatlabFile[opciones]{ruta/al/archivo.m}
\newcommand{\MatlabFile}[2][]{%
  \lstinputlisting[style=myMatlab,#1]{#2}%
}

% (Opcional) Soporte para acentos/símbolos en comentarios
\lstset{
  literate=
    {á}{{\'a}}1 {é}{{\'e}}1 {í}{{\'\i}}1 {ó}{{\'o}}1 {ú}{{\'u}}1
    {Á}{{\'A}}1 {É}{{\'E}}1 {Í}{{\'I}}1 {Ó}{{\'O}}1 {Ú}{{\'U}}1
    {ñ}{{\~n}}1 {Ñ}{{\~N}}1
    {μ}{{$\mu$}}1 {Δ}{{$\Delta$}}1 {°}{{$^\circ$}}1
}